\section{Outline of the argument}
\label{ex-outline}

The construction starts with finite-dimensional Wilson integrals and keeps
one choice of lattice observables through renormalization, extraction and
reconstruction. This section explains the route and the role of the main
estimates. It is a guide to the proofs and their hypotheses, not a
replacement for them.

\subsection{The construction and its scope}
\label{ex-outline-status}

The main statement is Theorem~\ref{thm:ex-construction-headline}.
Its assembly in Theorem~\ref{thm:ex-construction-assembly} retains
one lattice family, one refinement, one field quotient and one normalized
stress tensor. The weaker public target in Section~\ref{ex-statement}
is a consequence, but its abstract formulation alone does not identify
the lattice representatives with curvature polynomials. That connection
is supplied by the local sections and their reference-normalized transport.
The final assembly identifies the actual suppliers of the construction
premises, rather than inferring them from the existence of separate local
lemmas.

The result has no mass-gap or uniqueness-of-the-infinite-volume-limit
assertion. Its trace law excludes a specified class of finite mass-shell
profiles; it does not exclude every generalized free-field description.
The scope paragraph following the headline compares these conclusions
with the Yang--Mills existence and mass-gap problem.

\subsection{Prerequisites}
\label{ex-outline-prerequisites}

We assume finite-dimensional Lie groups, probability Haar measure on compact
groups, and measure and integration, including Fubini--Tonelli, the
$\pi$--$\lambda$ theorem and Bochner integration. We also use Cauchy's
theorem and formula, Hahn--Banach, elementary Hilbert-space theory, the
Fourier transform on Schwartz space, and elementary Gaussian inequalities
and Wick's rule. The specific uses of commutative $C^*$-algebras,
Riesz representation, the spectral theorem for unitary representations,
and the local analytic structure of matrix groups are cited where they
enter. In the spectral arguments the dual of $\mathbb R^d$ is identified
with $\mathbb R^d$ by $p\mapsto(x\mapsto e^{ip\cdot x})$; Fourier
normalizations are stated with the formulas that use them.

\subsection{Context}
\label{ex-outline-context}

The finite-volume measures are Wilson's lattice gauge theory
\cite{Wilson74}. Reflection positivity of these measures and the
resulting positive transfer matrix were established by Osterwalder and
Seiler and by L\"uscher \cite{OsterwalderSeiler,Luscher77}, and lattice
gauge theories were analysed from the constructive point of view by
Brydges, Fr\"ohlich and Seiler and by Seiler
\cite{BFS79,BFS80,BFS81,Seiler82}. The perturbative discovery of
asymptotic freedom \cite{GrossWilczek73,Politzer73} motivates the
logarithmic running law. The relation between the bare lattice coupling
and continuum renormalization schemes was computed at one loop in
\cite{HasenfratzHasenfratz80,DashenGross81} and at two loops in
\cite{LuscherWeiszCoupling}. None of these results supplies the uniform
lattice estimate proved here.

The construction proceeds by block-spin renormalization for lattice gauge
fields, an approach developed in Ba\l{}aban's programme. Its ingredients
include regularity and decay of lattice Green's functions
\cite{Balaban83Green}; propagators for the renormalization transformations,
first for the free gauge field and then in a background field
\cite{Balaban84Prop1,Balaban84Prop2,Balaban85Background}; averaging
operations, spaces of regular configurations and gauge fixing
\cite{Balaban85Averaging,Balaban85Spaces}; the variational problem for
background fields \cite{Balaban85Variational}; ultraviolet stability in
three dimensions \cite{Balaban85UV3}; effective actions, cluster
expansions and coupling-constant renormalization in four dimensions
\cite{Balaban87RG1,Balaban88RG2,Balaban88Convergent}; and the large-field
renormalization operation \cite{Balaban89LF1,Balaban89LF2}. Dimock's
expository series \cite{Dimock13I,Dimock13II,Dimock13III} presents the
method in detail for a scalar model. Other constructive approaches to
non-abelian gauge fields are Federbush's phase-cell programme
\cite{Federbush86I,Federbush87III,Federbush88IV,Federbush90V,Federbush87VI}
and the construction of Magnen, Rivasseau and S\'en\'eor with an
infrared cutoff \cite{MRS93}. Continuum limits are known for abelian
lattice gauge theories in three and four dimensions
\cite{Gross83,Driver87} and for the abelian Higgs model \cite{King86},
and the two-dimensional Yang--Mills measure has been constructed in the
continuum \cite{Driver89,Sengupta97,Levy03}. These papers provide
historical context; their construction theorems are not used as
substitutes for the estimates below.

For probabilistic background, see \cite{Chatterjee19,Chevyrev22}.
Subsequent work constructs the Langevin dynamics for the two-dimensional
Yang--Mills measure \cite{CCHS22}. In three dimensions, work constructs
state spaces for distributional gauge orbits \cite{CaoChatterjee24}
and a local Markov process on gauge orbits for the stochastic
Yang--Mills--Higgs flow \cite{CCHS24}.
The existence and mass-gap problem is formulated in
\cite{JaffeWitten06}.

The passage to a relativistic theory follows the Osterwalder--Schrader
strategy \cite{OS73,OS75}. The 1975 sequel corrects the original
sufficiency argument and gives revised hypotheses. Glaser
\cite{Glaser74} proves analytic continuation and a corresponding
positivity property on the causal domain. Zinoviev \cite{Zinoviev95}
proves an equivalence theorem with a further weak spectral condition. The
Wightman framework and the reconstruction of a relativistic theory from
Euclidean correlations are treated in \cite{StreaterWightman,GlimmJaffe87}.
Here the continuation,
boundary, locality and spectral arguments are carried out directly for
the constructed correlations. The comparison with the corrected abstract
theorem in Subsection~\ref{ex-reflection} is not a premise of that proof.

\subsection{The order of the limits}
\label{ex-outline-limits}

The lattice spacings $a_N=D_N^{-1}$, physical periods $P_p$, bare
coefficients and graded insertions are chosen before the requested
extraction order. For each pair of increasing index maps $(u,p)$ one
then chooses a refinement $(u\circ v,p\circ q)$ and one continuum
hierarchy. The Wilson limit uses the \emph{full} mesh sequence at every
suitable fixed period, followed by the selected period sequence $p\circ q$.
The refined mesh is used for auxiliary comparisons, not as a replacement
for that full mesh limit. One refinement serves all arities, labels and
tests. Lemma~\ref{ex-full-mesh} identifies a subsequential row limit
with an already established full row limit; it does not prove full mesh
convergence from compactness.

\subsection{The plan in ten steps}
\label{ex-outline-plan}

\begin{enumerate}
\item \emph{Finite lattice and quadratic transformation.}
Normalize the Wilson measure, establish gauge and reflection positivity,
and calculate the free block transformation. Laurent-symbol estimates give
exponentially decaying rows on every torus. The flat one-step estimate and
its scalar recurrence compare different terminal cutoffs. Curvature
factorization supplies the positive rows used by the quadratic horizon;
its bounds retain the powers of the block side in the residue sums.
See Sections~\ref{ex-lattice}--\ref{sec:finite-menu-history-handoff}.

\item \emph{One nonlinear block step.}
Defect geometry allocates action to large-field regions. On each sector,
a positive precision and localized covariance control the Gaussian
integral. Exact reassembly, a convergent cluster expansion and local jet
extraction give the nonlinear step and its marked versions. These are
identities and estimates for the same density, not separately chosen
approximations. The principal constructions begin in
Sections~\ref{sec:local-transfer}, \ref{sec:cluster-expansion}
and~\ref{sec:gaussian-reassembly}.

\item \emph{Stationary analysis and the one-loop coefficient.}
A uniform stationary chart and determinant comparison isolate the
quadratic logarithmic term. The annulus estimate gives
$\lambda_G=11C_G/(24\pi^2)$ in the color inner product fixed by the
Wilson representation. It is the same constant as $2b_0(G)$ and the
shooting slope $s_G$; see \eqref{eq:shoot-slope} and
Subsection~\ref{sec:running-coefficient}.

\item \emph{One selected history and schedule.}
The reset and shooting construction keeps the couplings within a
sublinear error of a line in $\log a_m^{-1}$. Specialization to the
finest coefficient gives the period-uniform running law. Numerical
profiles and floors are chosen in the required order before the block
lengths; costs evaluated after those choices are absorbed on the same
schedule. The public family retains the resulting actual coefficients,
levels, paths and insertion transport. See
Sections~\ref{sec:surviving-interval-shooting},
\ref{sec:common-scale-selection} and~\ref{sec:public-run-binding}.

\item \emph{Centered estimates and ordered limits.}
Reflection positivity turns finite-cell root bounds into estimates for
separated correlations. A fixed intermediate scale controls the discarded
higher marked terms; convergence of the retained data there yields the
full mesh Cauchy estimate. Difference estimates give smoothness away from
coincidences, and compactness gives one volume refinement for all tests.
See Sections~\ref{sec:fullmesh-correlation-cauchy},
\ref{sec:fixed-rung-convergence}, \ref{ex-smooth-limits}
and~\ref{ex-public-same-run-cauchy}. Their common-history premises are
part of the final assembly check.

\item \emph{The Euclidean Hilbert space and phases.}
Reflection positivity gives a contraction semigroup. Bounded slab
insertions compressed to its ground space generate a commutative algebra.
Its character space carries the vacuum's full-support probability
measure. Evaluating whole words gives component correlations and the
barycenter identity. Uniform slab bounds, derivatives and all-spectator
null relations survive this construction. See
Sections~\ref{ex-reflection}--\ref{ex-phase-fields}.

\item \emph{Fields, stress and calibrated actions.}
Finite probes and interacting/free comparisons identify field ranks and
compatible grade inclusions. The resulting physical field space carries
one conserved stress tensor with fixed Ward normalization, and one
compatible rotation and derivative calculus. Its shell identities pass
to every character before continuation. See
Sections~\ref{sec:finite-rank-kernel-transport},
\ref{sec:actual-stress-normalization},
\ref{sec:physical-angular-actions} and~\ref{sec:phase-ward-inheritance}.

\item \emph{Relativistic reconstruction in every character.}
Complex heat time gives ordered holomorphic correlations; common growth
bounds give joint tempered boundary distributions. Lorentz covariance,
locality, the word domain and the forward spectrum are obtained before
vacuum uniqueness. The latter first holds almost everywhere, then for
every character: locality and cone support give a uniform low-energy
bound, hence a uniform decaying heat remainder; continuity in the
character and full support remove the exceptional set.
Theorems~\ref{thm:joint-boundary}, \ref{thm:field-spectrum}
and~\ref{thm:ex-every-character-vacuum} give these steps.
Proposition~\ref{prop:assembly-boundary-barycenter} retains the mixture
identity on arbitrary joint Schwartz tests.

\item \emph{The operator-product expansion.}
One finite-dimensional prescription is matched with the local grouping
of inserted fields. Uniform character estimates give both coefficient
bounds and the remainder after truncation at any fixed degree. The
coefficients are independent of the state and spectators, whereas the
small-radius threshold may depend on the fields, geometry and truncation.
See Theorem~\ref{thm:ope-quantitative} and the matching arguments in
Subsection~\ref{sec:reference-match}.

\item \emph{The centered physical trace.}
The same Ward-normalized tensor is centered at the selected character.
Rounded lattice comparisons, the common chart estimates and the
free contraction determine its short-distance amplitude. Positivity and
monotonicity remove the rounding; the finite-profile argument then gives
the stated exclusion. Theorem~\ref{thm:trace-selected-endpoint}
attaches this conclusion to the selected hierarchy without a new
extraction or a new stress normalization.
\end{enumerate}

\subsection{Reading order and final interfaces}
\label{ex-outline-map}
\label{ex-outline-open}

The exposition presents the analytic consequences before many of their
construction estimates. Forward references to a later producer retain
that producer's hypotheses; they do not identify two independently
chosen families. The index in Section~\ref{sec:symbolic-index} records
the principal symbols and normalization changes.

The final assembly in Subsection~\ref{ex-assembly} follows the order:
selected finite family, full mesh and volume limits, common field
quotients, all-character reconstruction and OPE, then the selected physical
trace. The actual background estimates are supplied by the one-step and
retained path construction of Subsection~\ref{sec:path-attachment-producer}.
Proposition~\ref{prop:fullmesh-selected-supplier} supplies the uniform
all-period bounds for the calibrated insertions themselves. The final
assembly verifies these applications on the same data; the reference map
here is a guide to those proofs, rather than a substitute for them.
